\section{Exact zero order along a formal arc}
\label{sec:arc}

A generic formal arc on \(X\) is a morphism
\[
\gamma:\Spec F[[t]]\longrightarrow X
\]
over an extension field \(F/\C\), whose fraction point maps to the
generic point of \(X\). For a Cartier divisor \(D\), write
\(v_\gamma(D)\) for the order in \(t\) of its local equation pulled
back along the arc. Genericity ensures that a nonzero equation remains
nonzero.

\begin{corollary}
\label{cor:arc}
Let \(X\) be smooth
integral projective over \(\C\), let \(L=-K_X\) have a smooth
semipositive Hermitian metric, and let \(M\) be a line bundle such that
\(-M\) is pseudoeffective. Let \(\gamma\) be a generic formal arc on
\(X\). Suppose that effective integral divisors satisfy
\[
N_j\sim M+m_jL,\qquad
m_j\in\Z_{>0},\quad m_j\longrightarrow\infty,\qquad
v_\gamma(N_j)=o(m_j).
\]
Then, for some integer \(r>0\), there is an effective integral divisor
\[
D\sim rL,\qquad v_\gamma(D)=0.
\]
\end{corollary}

\begin{proof}
If \(X\) is a point, use its zero divisor. Otherwise set \(n=\dim X>0\)
and choose \(G_0\) and \(\ell\) as in
Theorem~\ref{V:exact}. Smooth semipositivity makes \(L\)
nef. If \(\nu<n\), the functional \(\ell\) is nonnegative on effective
divisors, since its remaining factors are nef; by continuity it is
nonnegative on the closed pseudoeffective cone. If \(\alpha\) is the
curvature form normalized to represent \(c_1(L)\), then
\(\alpha^{\nu+1}=0\), so \(\ell(L)=0\). Thus
\[
0\leq\ell(N_j)
=\ell(M)+m_j\ell(L)
=\ell(M)\leq0.
\]
The last inequality uses pseudoeffectivity of \(-M\). This proves
\(\ell(N_j)=0\). For \(\nu=n\), the same conclusion follows from
the definition \(\ell=0\).

The generic fraction point of the arc gives an injection
\[
\C(X)\hookrightarrow F((t)).
\]
Restrict \(\ord_t\) to \(\C(X)^\times\). Its image is either zero
or \(d\Z\) for an integer \(d>0\). Let \(\xi\) be the image of the
closed point of the arc. The local ring \(\OO_{X,\xi}\) maps into
\(F[[t]]\), its maximal ideal maps into \(tF[[t]]\), and its units
have order zero. Hence the restricted valuation is centered at \(\xi\).
If it is zero, this forces \(\xi\) to be the generic point; we use
the zero-valuation case of Theorem~\ref{V:exact}. Otherwise divide by
\(d\) to obtain a normalized discrete valuation centered at \(\xi\).
Orders of local divisor equations agree with arc orders, up to this
fixed normalization.

The assumption \(v_\gamma(N_j)=o(m_j)\) is exactly the required
normalized-order limit, and dividing by \(d\) leaves that limit unchanged.
Theorem~\ref{V:exact} therefore supplies an effective
rational divisor \(N\sim_{\Q}L\) with \(v_\gamma(N)=0\). Choose \(r>0\)
clearing the coefficients of \(N\) and the rational linear equivalence.
Then \(D=rN\) is effective integral, \(D\sim rL\), and
\(v_\gamma(D)=0\).
\end{proof}

In particular, bounded input orders imply the corollary, and the same
output \(D\) satisfies \(v_\gamma(D)<\delta r\) for every \(\delta>0\).
There is also an independent route to the weaker approximate conclusion.
It starts from the numerical nef reduction of Bauer et al.\ 
\cite[Theorem~2.1]{BCEKPRSW2002} and uses the multiplier-ideal quotient
extension theorem of Cao--Demailly--Matsumura
\cite[Theorem~1.1]{CDM2017} to obtain the required jets. The passage
from nef reduction to linear equivalence, and the construction of that
quotient, are additional steps; the corresponding anticanonical
argument assumes irregularity zero. The implication proved here uses
neither that assumption nor the extension theorem. It uses the stronger
exact valuation theorem, whose finite-generation preparation remains essential.

